% 原章第22—28页：6.2.1.4至式6.97。
\subsubsection{卷积神经网络}

在自编码器中，一层的每个神经元都与前一层的全部神经元连接，待估计的参数数目会非常大。例如，假定一幅图像有 $100\times100$ 个像素，则与输入层神经元连接的神经元数目应为 10,000。因此，有必要限制神经元之间的连接。卷积神经网络（convolutional neural network，CNN）通过局部连接（local connectivity）实现这一目标（LeCun 1989）。例如，图 6.4 中，CNN 仅将每个隐藏神经元与输入向量中少数相邻神经元连接。深度 CNN 在多层中施加这种局部连接。
\begin{figure}[htbp]
\centering
\begin{tikzpicture}[x=1.15cm,y=1cm,>=Stealth]
\node[anchor=east] at (-.25,1.5) {隐藏层};\node[anchor=east] at (-.25,0) {输入层};
\foreach \k in {1,...,6} \node[circle,draw=blue!60!black,fill=blue!60,minimum size=5mm] (chSixF4i\k) at (\k-1,0) {};
\foreach \k in {1,...,4} {
 \node[circle,draw=blue!60!black,fill=blue!60,minimum size=5mm] (chSixF4h\k) at (\k,1.5) {};
 \pgfmathtruncatemacro{\kp}{\k+1}\pgfmathtruncatemacro{\kpp}{\k+2}
 \draw[->,red] (chSixF4i\k)--(chSixF4h\k);
 \draw[->,green!65!black] (chSixF4i\kp)--(chSixF4h\k);
 \draw[->,blue!65!black] (chSixF4i\kpp)--(chSixF4h\k);
}
\end{tikzpicture}
\origfig{4}{局部连接神经元示意图。}
\end{figure}


CNN 源于神经网络在影像数据分析中的应用，旨在处理多重数组数据，并利用图像像素之间的空间信息。CNN 的体系结构由一系列不同的层堆叠而成，将输入体变换为输出体：卷积层、池化层、修正线性单元（rectified linear unit，ReLU）层、池化层、全连接层以及损失层。卷积是 CNN 的核心组成部分，下面首先研究卷积。

\bookdeepheading{paragraph}{6.2.1.4.1}{卷积或互相关运算}

卷积是一种滤波器，可从数据中提取特征、去除噪声。设 $f(t)$ 为含噪信号函数。为提高精度、得到噪声较小的信号，对时刻 $t$ 附近的若干次测量取平均。越近的测量越相关，因此，对邻近测量取平均时，应给较近的测量更大权重。令 $g(t)$ 为权函数，通常称为滤波器（filter）或核（kernel）。以下称为卷积的加权平均运算，可估计噪声较小的信号：
\begin{equation}
s(t)=(f*g)(t)=\int_{-\infty}^{\infty}f(\tau)g(t-\tau)\,d\tau,
\tag{6.85}
\end{equation}
其中，$*$ 表示卷积。第一个参数称为输入，第二个参数称为核，输出通常称为特征图（feature map）。

实际中，只能按固定间隔测量信号。假定信号 $f(t)$ 与核 $g(t)$ 只能在整数 $t$ 处取值，则式 (6.85) 中的积分可离散化，离散卷积定义为
\begin{equation}
s(t)=(f*g)(t)=\sum_{-\infty}^{\infty}f(\tau)g(t-\tau).
\tag{6.86}
\end{equation}
在影像数据分析中，使用二维图像 $A$ 作为输入，因而需要将一维卷积推广为二维卷积。设图像信号（或激活量）为
\[
A:\{a_{i,j},\ i=1,\ldots,n_1,\ j=1,\ldots,n_2\},
\]
核为
\[
W=\{w_{u,v},\ u=0,\ldots,h_1,\ v=0,\ldots,h_2\}.
\]
图像或激活量 $A$ 与核 $W$ 的离散卷积定义为
\begin{equation}
(A*W)_{ij}=\sum_{u=0}^{h_1}\sum_{v=0}^{h_2}w_{u,v}a_{i-u,j-v}.
\tag{6.87}
\end{equation}
令 $m=i-u$、$n=j-v$，则式 (6.86) 变为
\begin{equation}
(A*W)_{ij}=\sum_m\sum_na_{m,n}w_{i-m,j-n}.
\tag{6.88}
\end{equation}
实现时，常用互相关（cross-correlation）代替卷积运算，其定义为
\begin{equation}
(A*W)_{ij}=\sum_m\sum_na_{i+m,j+n}w_{m,n}.
\tag{6.89}
\end{equation}

\Needspace{4\baselineskip}
\noindent\textbf{例 6.1}

考虑一个 $4\times4$ 的输入矩阵 $A$ 和一个 $2\times2$ 的核矩阵 $W$（图 6.5）。$W$ 与 $A$ 的卷积生成一个 $3\times3$ 的输出矩阵，即特征图。
\begin{figure}[htbp]
\centering
\begin{tikzpicture}[x=1cm,y=1cm,>=Stealth]
\filldraw[fill=red!75,draw=white] (0.0,2.4) rectangle ++(0.9,-0.38);
\node[font=\small] at (0.45,2.21) {1};
\filldraw[fill=red!75,draw=white] (0.9,2.4) rectangle ++(0.9,-0.38);
\node[font=\small] at (1.35,2.21) {0};
\filldraw[fill=blue!60,draw=white] (1.8,2.4) rectangle ++(0.9,-0.38);
\node[font=\small] at (2.25,2.21) {1};
\filldraw[fill=blue!60,draw=white] (2.7,2.4) rectangle ++(0.9,-0.38);
\node[font=\small] at (3.15,2.21) {0};
\filldraw[fill=red!75,draw=white] (0.0,2.02) rectangle ++(0.9,-0.38);
\node[font=\small] at (0.45,1.8299999999999998) {0};
\filldraw[fill=red!75,draw=white] (0.9,2.02) rectangle ++(0.9,-0.38);
\node[font=\small] at (1.35,1.8299999999999998) {2};
\filldraw[fill=blue!15,draw=white] (1.8,2.02) rectangle ++(0.9,-0.38);
\node[font=\small] at (2.25,1.8299999999999998) {0};
\filldraw[fill=blue!15,draw=white] (2.7,2.02) rectangle ++(0.9,-0.38);
\node[font=\small] at (3.15,1.8299999999999998) {1};
\filldraw[fill=blue!10,draw=white] (0.0,1.64) rectangle ++(0.9,-0.38);
\node[font=\small] at (0.45,1.45) {1};
\filldraw[fill=blue!10,draw=white] (0.9,1.64) rectangle ++(0.9,-0.38);
\node[font=\small] at (1.35,1.45) {2};
\filldraw[fill=blue!10,draw=white] (1.8,1.64) rectangle ++(0.9,-0.38);
\node[font=\small] at (2.25,1.45) {1};
\filldraw[fill=blue!10,draw=white] (2.7,1.64) rectangle ++(0.9,-0.38);
\node[font=\small] at (3.15,1.45) {1};
\filldraw[fill=blue!15,draw=white] (0.0,1.2599999999999998) rectangle ++(0.9,-0.38);
\node[font=\small] at (0.45,1.0699999999999998) {1};
\filldraw[fill=blue!15,draw=white] (0.9,1.2599999999999998) rectangle ++(0.9,-0.38);
\node[font=\small] at (1.35,1.0699999999999998) {3};
\filldraw[fill=blue!15,draw=white] (1.8,1.2599999999999998) rectangle ++(0.9,-0.38);
\node[font=\small] at (2.25,1.0699999999999998) {0};
\filldraw[fill=blue!15,draw=white] (2.7,1.2599999999999998) rectangle ++(0.9,-0.38);
\node[font=\small] at (3.15,1.0699999999999998) {2};\filldraw[fill=red!75,draw=white] (4.7,2.0) rectangle ++(0.9,-0.38);
\node[font=\small] at (5.15,1.81) {1};
\filldraw[fill=red!75,draw=white] (5.6000000000000005,2.0) rectangle ++(0.9,-0.38);
\node[font=\small] at (6.050000000000001,1.81) {0};
\filldraw[fill=red!75,draw=white] (4.7,1.62) rectangle ++(0.9,-0.38);
\node[font=\small] at (5.15,1.43) {0};
\filldraw[fill=red!75,draw=white] (5.6000000000000005,1.62) rectangle ++(0.9,-0.38);
\node[font=\small] at (6.050000000000001,1.43) {1};\filldraw[fill=yellow!80!orange,draw=white] (1.4,-0.1) rectangle ++(0.9,-0.38);
\node[font=\small] at (1.8499999999999999,-0.29000000000000004) {3};
\filldraw[fill=yellow!80!orange,draw=white] (2.3,-0.1) rectangle ++(0.9,-0.38);
\node[font=\small] at (2.75,-0.29000000000000004) {0};
\filldraw[fill=yellow!80!orange,draw=white] (3.2,-0.1) rectangle ++(0.9,-0.38);
\node[font=\small] at (3.65,-0.29000000000000004) {2};
\filldraw[fill=blue!15,draw=white] (1.4,-0.48) rectangle ++(0.9,-0.38);
\node[font=\small] at (1.8499999999999999,-0.67) {2};
\filldraw[fill=blue!15,draw=white] (2.3,-0.48) rectangle ++(0.9,-0.38);
\node[font=\small] at (2.75,-0.67) {3};
\filldraw[fill=blue!15,draw=white] (3.2,-0.48) rectangle ++(0.9,-0.38);
\node[font=\small] at (3.65,-0.67) {1};
\filldraw[fill=blue!10,draw=white] (1.4,-0.86) rectangle ++(0.9,-0.38);
\node[font=\small] at (1.8499999999999999,-1.05) {4};
\filldraw[fill=blue!10,draw=white] (2.3,-0.86) rectangle ++(0.9,-0.38);
\node[font=\small] at (2.75,-1.05) {2};
\filldraw[fill=blue!10,draw=white] (3.2,-0.86) rectangle ++(0.9,-0.38);
\node[font=\small] at (3.65,-1.05) {3};
\node[anchor=west] at (0,.55) {输入 $A$};\node at (5.6,2.3) {核 $W$};
\node[anchor=west] at (4.2,-.5) {特征图 $(A*W)$};
\draw[->,green!65!black] (0,2.02)--(-.3,2.02)--(-.3,-.3)--(1.4,-.3);
\draw[->,green!65!black] (5.6,1.24)--(5.6,.3)--(1.5,.3)--(1.5,-.1);
\end{tikzpicture}
\origfig{5}{卷积示例。}
\end{figure}


CNN 的结构见图 6.6。CNN 包含一组可学习的核，每个核都是具有宽、高、深度的三维矩阵。例如，尺寸 $6\times6\times3$ 表示宽和高均为 6，以及输入数据的深度。沿前一层数据的宽、高方向滑动核，可生成二维激活图或特征图；它输出核在前一层各位置的响应。
\begin{figure}[htbp]
\centering
\begin{tikzpicture}[x=1cm,y=1cm,>=Stealth,draw=blue!60!black]
\draw (0,0) rectangle (.2,2.8);\draw (.2,0)--(1.2,1.3)--(1.2,4.1)--(.2,2.8);
\draw (0,2.8)--(1,4.1)--(1.2,4.1);\node at (.1,-.25) {$3$};
\node at (1.45,3.8) {$32$};\node at (1,1) {$32$};
\draw (.45,1.6)--(.65,1.4)--(.65,2.2)--(.9,2.5)--(.65,2.6)--(.45,2.35)--cycle;
\node[black] at (.85,2.9) {$A$};
\foreach \xx/\yy in {1.8/.6,2.05/.35,2.3/.1} \draw[fill=white] (\xx,\yy) rectangle ++(2,2.8);
\node[black] at (3.3,2.4) {特征图};
\filldraw[fill=blue!60] (2.9,1.2) rectangle (3.2,1.75);
\draw (.65,1.4)--(2.9,1.2);\draw (.65,2.2)--(2.9,1.75);\draw (.9,2.5)--(3.2,1.75);
\node[black] at (2.8,.7) {$W$};\draw[->] (2.8,.9)--(3,1.2);
\node[black,anchor=west] at (3.55,1.5) {核};\draw[->] (3.55,1.5)--(3.2,1.5);
\draw[<->] (4.8,.5)--(4.8,3.8);\draw (4.3,.1)--(5,.8);\draw (4.3,3.5)--(5,4.2);
\node[black,anchor=west] at (5.15,2.8) {深度：核的数目};
\end{tikzpicture}
\origfig{6}{卷积层的结构。}
\end{figure}


卷积层的深度定义为卷积运算所用核的数目。图 6.6 的网络使用三个核对原始图像进行卷积，因此生成三个不同的特征图。这三个特征图可看作堆叠的二维矩阵，特征图的“深度”为 3。深度列（depth column）定义为与输入同一区域进行卷积的一组神经元，常用于研究参数共享。

步长（stride）定义为在输入矩阵上滑动核时所移动的像素数。步长为 1，表示每次移动一个像素；步长为 2，则每次跳过两个像素。为控制输出大小，常在输入数据边界周围补零。加入零填充的卷积称为宽卷积（wide convolution），不加入零填充的卷积称为窄卷积（narrow convolution）。

设计卷积层时，需要根据输入体的尺寸 $W$、核的尺寸 $F$、核移动的步长 $S$，以及边界使用的零填充量 $P$，计算卷积层的空间尺寸。卷积层在一个维度上的尺寸为
\begin{equation}
m=\frac{W-F+2P}{S}+1.
\tag{6.90}
\end{equation}

\Needspace{4\baselineskip}
\noindent\textbf{例 6.2}

考虑图 6.7 所示卷积层的空间排列。假定 $W=7$、$F=3$、$P=1$、$S=2$，则输出（特征图）大小为 $m=(7-3+2)/2+1=4$。特征图神经元内部的数字，是输入与核向量 $[1,0,-1]$ 的卷积值。核向量的元素就是权重，也用红、绿、紫三种颜色表示。这些权重由特征图中的四个神经元共享。
\begin{figure}[htbp]
\centering
\begin{tikzpicture}[x=.85cm,y=.8cm,>=Stealth,
chSixFsevenN/.style={draw=blue!60!black,fill=blue!60,text=white,minimum width=5mm,minimum height=4mm,inner sep=1pt,font=\small}]
\foreach \xx/\vv in {0/0,1/3,2/2,3/1,4/0,5/1,6/2,7/3,8/0}
 \node[chSixFsevenN] (chSixF7i\xx) at (\xx,0) {$\vv$};
\foreach \xx/\vv/\kk in {1/-2/0,3/2/2,5/-2/4,7/2/6} {
 \node[chSixFsevenN,fill=cyan!70!blue] (chSixF7h\xx) at (\xx,2) {$\vv$};
 \pgfmathtruncatemacro{\kp}{\kk+1}\pgfmathtruncatemacro{\kpp}{\kk+2}
 \draw[->,red] (chSixF7h\xx)--(chSixF7i\kk);
 \draw[->,green!70!black] (chSixF7h\xx)--(chSixF7i\kp);
 \draw[->,violet] (chSixF7h\xx)--(chSixF7i\kpp);
}
\draw[red,dash dot] (-.15,1.65) rectangle (7.95,2.65);\node[red] at (4,2.4) {特征图};
\node[font=\small] at (.45,1.15) {$1$};\node[font=\small] at (.82,.65) {$0$};\node[font=\small] at (1.5,1.35) {$-1$};
\draw[blue!60!black] (4.5,3.5) rectangle (8.4,4.2);
\foreach \xx/\vv in {5/1,6.5/0,8/-1} \node[chSixFsevenN] at (\xx,3.85) {$\vv$};
\node at (6.5,3.2) {核};
\end{tikzpicture}
\origfig{7}{卷积层的空间排列。}
\end{figure}


\Needspace{4\baselineskip}
\noindent\textbf{例 6.3}

考虑一个尺寸为 $227\times227\times3$ 的大型深度 CNN，用来对大型标注数据集 Image Net 中的 120 万幅图像进行分类（Krizhevsky et al. 2012）。卷积层的每个神经元使用尺寸为 $11\times11$ 的核（$F=11$），步长 $S=4$，不作零填充（$P=0$）。根据式 (6.89)，输出（特征图）大小为 $m=(227-11+0)/4+1=55$。卷积层使用了 $K=96$ 个不同的核（深度），因此包含 96 个特征图，每个特征图大小为 $55\times55$。卷积层总共使用 $55\times55\times96=290,400$ 个神经元。每个神经元与前一层中大小为 $[11\times11\times3]$ 的区域连接。每个深度列中的 96 个神经元，都连接前一层中大小为 $[11\times11\times3]$ 的同一区域，但权重不同。因此，每个神经元有 $11\times11\times3=363$ 个权重。卷积层权重总数为 $290400\times363=105705600$。显然，网络中的参数数目过多。

为控制卷积层参数数目，应采用参数共享策略，即让同一个参数用于多个神经元。一个特征图常称为深度切片（depth slice），本例有 96 个深度切片。每个深度切片中的核使用一组独有的权重（$11\times11\times3=363$）。总参数数目为（$363\times96=34848$ 个权重 $+96$ 个偏置）34,544。可见，参数共享大幅减少了卷积层的参数数目。

卷积层是 CNN 体系结构的基本模块，由许多特征图组成。每个特征图由输入层或前一层特征图与特定的核卷积得到。假定第 $l$ 层为卷积层，包含 $m_1^{(l)}$ 个大小为 $m_2^{(l)}\times m_3^{(l)}$ 的特征图。卷积层的输入包含 $m_1^{(l-1)}$ 个大小为 $m_2^{(l-1)}\times m_3^{(l-1)}$ 的特征图。当 $l=1$ 时，输入是一幅图像。令 $A_i^{(l)}$ 与 $B_i^{(l)}$ 分别为第 $l$ 层第 $i$ 个特征图的激活矩阵与偏置矩阵：
\[
A_i^{(l)}=\begin{bmatrix}
a_{11}^{(l)}&\cdots&a_{1m_3^{(l)}}^{(l)}\\
\vdots&\vdots&\vdots\\
a_{m_2^{(l)}1}^{(l)}&\cdots&a_{m_2^{(l)}m_3^{(l)}}^{(l)}
\end{bmatrix},\qquad
B_i^{(l)}=\begin{bmatrix}
b_{11}^{(l)}&\cdots&b_{1m_3^{(l)}}^{(l)}\\
\vdots&\vdots&\vdots\\
b_{m_2^{(l)}1}^{(l)}&\cdots&b_{m_2^{(l)}m_3^{(l)}}^{(l)}
\end{bmatrix}.
\]
将第 $l-1$ 层第 $j$ 个特征图连接到第 $l$ 层第 $i$ 个特征图的核矩阵记为 $W_{i,j}^{(l)}$，其大小为 $(2h_1^{(l)}+1)\times(2h_2^{(l)}+1)$，定义如下：
\[
W_{i,j}^{(l)}=\begin{bmatrix}
(W_{i,j}^{(l)})_{-h_1^{(l)},-h_2^{(l)}}&\cdots&(W_{i,j}^{(l)})_{-h_1^{(l)},h_2^{(l)}}\\
\vdots&\vdots&\vdots\\
(W_{i,j}^{(l)})_{h_1^{(l)},-h_2^{(l)}}&\cdots&(W_{i,j}^{(l)})_{h_1^{(l)},h_2^{(l)}}
\end{bmatrix}.
\]
通过卷积计算激活矩阵 $A_i^{(l)}$：
\begin{equation}
\begin{aligned}
(A_i^{(l)})_{r,s}&=(B_i^{(l)})_{r,s}\\
&\quad+\sum_{j=1}^{m_1^{(l-1)}}\sum_{u=-h_1^{(l)}}^{h_1^{(l)}}\sum_{v=-h_2^{(l)}}^{h_2^{(l)}}
(W_{i,j}^{(l)})_{u,v}(A_j^{(l-1)})_{r+u,s+v}.
\end{aligned}\tag{6.91}
\end{equation}

\bookdeepheading{paragraph}{6.2.1.4.2}{非线性层（ReLU）}

大部分真实数据是非线性的，而卷积是线性运算，因此每次卷积后通常使用非线性映射，捕捉原始数据的非线性。假定第 $l$ 层为非线性层，第 $l-1$ 层为其输入层。令 $m_1^{(l)}$ 与 $m_1^{(l-1)}$ 分别表示两层的特征图数目，并假定各特征图大小相同，即 $m_2^{(l-1)}=m_2^{(l)}$、$m_3^{(l-1)}=m_2^{(l-1)}$。令 $A_i^{(l)}$ 与 $A_i^{(l-1)}$ 分别为非线性层（第 $l$ 层）与输入层（第 $l-1$ 层）中第 $i$ 个特征图的激活矩阵，非线性映射为
\begin{equation}
(A_i^{(l)})_{rs}=\sigma((A_i^{(l-1)})_{rs}),
\tag{6.92}
\end{equation}
其中，$\sigma$ 为非线性激活函数，可以取双曲正切函数 $\tanh(x)=(e^x-e^{-x})/(e^x+e^x)$，或 Sigmoid 函数 $\operatorname{sigm}(x)=1/(1+e^{-x})$。

然而，Sigmoid 函数往往有严重局限。在神经网络的反向传播过程中，使用 Sigmoid 函数激活时，误差在每一层都会被压缩（至少）四分之一。因此，网络越深，数据中“丢失”的信息就越多；输出层即使存在一些较大误差，也可能无法影响前面各层神经元的权重。

为克服 Sigmoid 函数作为非线性映射函数的局限，近来的 CNN 常使用修正线性单元（ReLU）执行非线性运算。假定第 $l$ 层为修正线性层，ReLU 中的激活定义为
\begin{equation}
(A_i^{(l)})_{r,s}=\max\bigl(0,(A_i^{(l-1)})_{r,s}\bigr).
\tag{6.93}
\end{equation}
ReLU 是最简单的非线性激活函数。它作用于前一层的每个神经元，属于逐元素运算。ReLU 运算后，来自特征图的全部负输入信号均被替换为零。执行非线性映射的阶段通常称为检测阶段（detector stage）。

\bookdeepheading{paragraph}{6.2.1.4.3}{特征池化与下采样层}

为减小表示的空间尺寸，减少特征图网络中的参数数目与计算量，常在 CNN 相邻卷积层之间周期性插入池化层（pooling layer）。空间池化的目的，是降低每个特征图的维数，同时保留最基本的信息。池化运算常用最大值、平均值和求和三类函数。

对输入特征图的每个深度切片，独立执行最大值运算，调整其空间尺寸。设第 $l$ 层为池化层，假定池化层与输入层的特征图数目相同，即 $m_1^{(l)}=m_1^{(l-1)}$。定义一个空间邻域，例如大小为 $u\times v$ 的窗口。最大池化在每个窗口中，取空间邻域内的最大元素。具体为
\begin{equation}
(A_i^{(l)})_{r,s}=\max_{u,v}\bigl((A_i^{(l-1)})_{r+u,s+v},\quad1\leq u\leq p,\ 1\leq v\leq q.
\tag{6.94}
\end{equation}
类似地，平均池化定义为
\begin{equation}
(A_i^{(l)})_{r,s}=\frac1{pq}\sum_{u=1}^p\sum_{v=1}^q(A_i^{(l-1)})_{r+u,s+v},
\tag{6.95}
\end{equation}
求和池化定义为
\begin{equation}
(A_i^{(l)})_{r,s}=\sum_{u=1}^p\sum_{v=1}^q(A_i^{(l-1)})_{r+u,s+v}.
\tag{6.96}
\end{equation}

\Needspace{4\baselineskip}
\noindent\textbf{例 6.4}

考虑一个大小为 $6\times6$ 的修正后的特征图与一个 $3\times3$ 的核（图 6.8）。令 $S=3$，用大小为 3、步长为 3 的核，对 $6\times6$ 输入矩阵进行池化，得到 $2\times2$ 输出矩阵。每次最大值运算在 9 个数（$3\times3$ 彩色方块）中进行。对红、绿、黄、紫四个彩色方块执行四次最大池化，生成 $2\times2$ 输出矩阵。
\begin{figure}[htbp]
\centering
\begin{tikzpicture}[x=1cm,y=1cm,>=Stealth]
\filldraw[fill=red!80,draw=white] (0.0,2.5) rectangle ++(0.45,-0.35);
\node[font=\small] at (0.225,2.325) {5};
\filldraw[fill=red!80,draw=white] (0.45,2.5) rectangle ++(0.45,-0.35);
\node[font=\small] at (0.675,2.325) {3};
\filldraw[fill=red!80,draw=white] (0.9,2.5) rectangle ++(0.45,-0.35);
\node[font=\small] at (1.125,2.325) {2};
\filldraw[fill=green!70!black,draw=white] (1.35,2.5) rectangle ++(0.45,-0.35);
\node[font=\small] at (1.575,2.325) {1};
\filldraw[fill=green!70!black,draw=white] (1.8,2.5) rectangle ++(0.45,-0.35);
\node[font=\small] at (2.025,2.325) {1};
\filldraw[fill=green!70!black,draw=white] (2.25,2.5) rectangle ++(0.45,-0.35);
\node[font=\small] at (2.475,2.325) {7};
\filldraw[fill=red!80,draw=white] (0.0,2.15) rectangle ++(0.45,-0.35);
\node[font=\small] at (0.225,1.975) {4};
\filldraw[fill=red!80,draw=white] (0.45,2.15) rectangle ++(0.45,-0.35);
\node[font=\small] at (0.675,1.975) {2};
\filldraw[fill=red!80,draw=white] (0.9,2.15) rectangle ++(0.45,-0.35);
\node[font=\small] at (1.125,1.975) {1};
\filldraw[fill=green!70!black,draw=white] (1.35,2.15) rectangle ++(0.45,-0.35);
\node[font=\small] at (1.575,1.975) {5};
\filldraw[fill=green!70!black,draw=white] (1.8,2.15) rectangle ++(0.45,-0.35);
\node[font=\small] at (2.025,1.975) {1};
\filldraw[fill=green!70!black,draw=white] (2.25,2.15) rectangle ++(0.45,-0.35);
\node[font=\small] at (2.475,1.975) {4};
\filldraw[fill=red!80,draw=white] (0.0,1.8) rectangle ++(0.45,-0.35);
\node[font=\small] at (0.225,1.625) {2};
\filldraw[fill=red!80,draw=white] (0.45,1.8) rectangle ++(0.45,-0.35);
\node[font=\small] at (0.675,1.625) {4};
\filldraw[fill=red!80,draw=white] (0.9,1.8) rectangle ++(0.45,-0.35);
\node[font=\small] at (1.125,1.625) {3};
\filldraw[fill=green!70!black,draw=white] (1.35,1.8) rectangle ++(0.45,-0.35);
\node[font=\small] at (1.575,1.625) {2};
\filldraw[fill=green!70!black,draw=white] (1.8,1.8) rectangle ++(0.45,-0.35);
\node[font=\small] at (2.025,1.625) {3};
\filldraw[fill=green!70!black,draw=white] (2.25,1.8) rectangle ++(0.45,-0.35);
\node[font=\small] at (2.475,1.625) {5};
\filldraw[fill=yellow!85,draw=white] (0.0,1.4500000000000002) rectangle ++(0.45,-0.35);
\node[font=\small] at (0.225,1.2750000000000001) {1};
\filldraw[fill=yellow!85,draw=white] (0.45,1.4500000000000002) rectangle ++(0.45,-0.35);
\node[font=\small] at (0.675,1.2750000000000001) {3};
\filldraw[fill=yellow!85,draw=white] (0.9,1.4500000000000002) rectangle ++(0.45,-0.35);
\node[font=\small] at (1.125,1.2750000000000001) {5};
\filldraw[fill=violet!70,draw=white] (1.35,1.4500000000000002) rectangle ++(0.45,-0.35);
\node[font=\small] at (1.575,1.2750000000000001) {8};
\filldraw[fill=violet!70,draw=white] (1.8,1.4500000000000002) rectangle ++(0.45,-0.35);
\node[font=\small] at (2.025,1.2750000000000001) {4};
\filldraw[fill=violet!70,draw=white] (2.25,1.4500000000000002) rectangle ++(0.45,-0.35);
\node[font=\small] at (2.475,1.2750000000000001) {2};
\filldraw[fill=yellow!85,draw=white] (0.0,1.1) rectangle ++(0.45,-0.35);
\node[font=\small] at (0.225,0.925) {6};
\filldraw[fill=yellow!85,draw=white] (0.45,1.1) rectangle ++(0.45,-0.35);
\node[font=\small] at (0.675,0.925) {2};
\filldraw[fill=yellow!85,draw=white] (0.9,1.1) rectangle ++(0.45,-0.35);
\node[font=\small] at (1.125,0.925) {3};
\filldraw[fill=violet!70,draw=white] (1.35,1.1) rectangle ++(0.45,-0.35);
\node[font=\small] at (1.575,0.925) {2};
\filldraw[fill=violet!70,draw=white] (1.8,1.1) rectangle ++(0.45,-0.35);
\node[font=\small] at (2.025,0.925) {1};
\filldraw[fill=violet!70,draw=white] (2.25,1.1) rectangle ++(0.45,-0.35);
\node[font=\small] at (2.475,0.925) {5};
\filldraw[fill=yellow!85,draw=white] (0.0,0.75) rectangle ++(0.45,-0.35);
\node[font=\small] at (0.225,0.5750000000000002) {3};
\filldraw[fill=yellow!85,draw=white] (0.45,0.75) rectangle ++(0.45,-0.35);
\node[font=\small] at (0.675,0.5750000000000002) {4};
\filldraw[fill=yellow!85,draw=white] (0.9,0.75) rectangle ++(0.45,-0.35);
\node[font=\small] at (1.125,0.5750000000000002) {2};
\filldraw[fill=violet!70,draw=white] (1.35,0.75) rectangle ++(0.45,-0.35);
\node[font=\small] at (1.575,0.5750000000000002) {3};
\filldraw[fill=violet!70,draw=white] (1.8,0.75) rectangle ++(0.45,-0.35);
\node[font=\small] at (2.025,0.5750000000000002) {5};
\filldraw[fill=violet!70,draw=white] (2.25,0.75) rectangle ++(0.45,-0.35);
\node[font=\small] at (2.475,0.5750000000000002) {3};\filldraw[fill=red!80,draw=white] (5.0,1.7) rectangle ++(0.55,-0.35);
\node[font=\small] at (5.275,1.525) {5};
\filldraw[fill=green!70!black,draw=white] (5.55,1.7) rectangle ++(0.55,-0.35);
\node[font=\small] at (5.825,1.525) {7};
\filldraw[fill=yellow!85,draw=white] (5.0,1.35) rectangle ++(0.55,-0.35);
\node[font=\small] at (5.275,1.175) {6};
\filldraw[fill=violet!70,draw=white] (5.55,1.35) rectangle ++(0.55,-0.35);
\node[font=\small] at (5.825,1.175) {8};
\node at (1.35,2.8) {$3\times3$ 核};\node at (1.35,.1) {步长 $S=3$};
\node at (5.55,2) {输出矩阵};\node at (3.75,1.85) {最大池化};
\draw[->,line width=3pt,green!60!black] (4.3,1.4)--(4.8,1.4);
\end{tikzpicture}
\origfig{8}{最大池化运算。}
\end{figure}


\bookdeepheading{paragraph}{6.2.1.4.4}{归一化层}

设第 $l$ 层为归一化层（normalization layer）。考虑减法归一化（subtractive normalization）与亮度归一化（brightness normalization）两种方法。先研究减法归一化。对于输入层 $l-1$ 中的每个特征图，它在归一化层 $l$ 中对应的输出激活定义为
\begin{equation}
A_i^{(l)}=A_i^{(l-1)}-\sum_{j=1}^{m_{l-1}}W_{G(\sigma)}*A_j^{(l-1)},
\tag{6.97}
\end{equation}
其中，高斯滤波器 $W_{G(\sigma)}$ 定义为
\[
(W_{G(\sigma)})_{r,s}=\frac1{\sqrt{2\pi\sigma^2}}\exp\left(\frac{r^2+s^2}{2\sigma^2}\right).
\]
